%! Piecewise-Defined Functions — Unit 2, Lesson 2.6 — Warm-Up (Answer Key)
\documentclass[10pt]{article}
\usepackage{atda-article}
\usepackage{atda-key}   % pulls in -boxes; do NOT also load -boxes
\newcommand{\TallMath}[1]{$\displaystyle #1\rule[-1.4em]{0pt}{3.2em}$}
\setlength{\workrowsep}{6pt}   % handwriting room; moves blank and key together

\begin{document}

\pageheader{Unit 2, Lesson 2.6}{Warm-Up --- Answer Key}
% NAMESTRIP: no \namedateperiod here — mirrors the blank. Teacher-only prose goes
% in the lesson plan as \begin{teachernote}[Warm-Up], never in a key.

\vspace{0.15in}

\begin{notesbox}{Two Rules, and a Pencil You Do Not Have to Lift}
\small
Five quick items. Item 2 is the one to slow down on --- read the policy twice before you compute.

\begin{enumerate}[leftmargin=1.6em, itemsep=4pt, topsep=4pt]

\item \textbf{Last night's homework.} Item 10 graphed a delivery fee: a flat \$5 for the first $3$
miles, then \$2 per mile after that. Give $F(2)$ and $F(6)$, and say how many \emph{rules} it took to
describe the whole graph.

\ansline{$F(2)=\$5$ and $F(6)=\$11$. It took \textbf{two} rules --- one flat, one climbing --- handing off at $m=3$.}\\

\item \textbf{A gym membership.} The plan charges \$25 a month for the first $6$ months and \$40 a
month for every month after that. \quad (a) What do you pay in month $3$? \quad (b) In month $9$?
\quad (c) What do you pay for a full $12$ months?

\ansline{(a) \$25 \quad (b) \$40 \quad (c) $6(25)+6(40)=150+240=\$390$ --- each rate covers only its own months.}\\

\item \textbf{Interval notation, both directions} (Lesson 2.1). \quad (a) Write $2 < w \le 5$ in
interval notation. \quad (b) Write $[0,40]$ as an inequality.

\ansline{(a) $(2,5]$ --- parenthesis excludes $2$, bracket includes $5$. \quad (b) $0 \le h \le 40$.}\\

\item \textbf{Can you trace it without lifting your pencil?} Answer yes or no for each graph.
\smallskip

\begin{center}
\begin{tabular}{@{}c@{\hspace{0.7cm}}c@{}}
\begin{tikzpicture}
\begin{axis}[
    width=5.4cm, height=3.6cm, title={\scriptsize (a)},
    xmin=0, xmax=6, ymin=0, ymax=8,
    xtick={0,1,2,3,4,5,6}, ytick={0,2,4,6,8},
    grid=both, grid style={linegray!45},
    axis lines=left, tick label style={font=\tiny}]
  \addplot[royal, very thick] coordinates {(0,2) (3,2) (6,8)};
\end{axis}
\end{tikzpicture}
&
\begin{tikzpicture}
\begin{axis}[
    width=5.4cm, height=3.6cm, title={\scriptsize (b)},
    xmin=0, xmax=6, ymin=0, ymax=8,
    xtick={0,1,2,3,4,5,6}, ytick={0,2,4,6,8},
    grid=both, grid style={linegray!45},
    axis lines=left, tick label style={font=\tiny}]
  \addplot[royal, very thick] coordinates {(0,2) (3,2)};
  \addplot[royal, very thick] coordinates {(3,6) (6,6)};
\end{axis}
\end{tikzpicture}
\end{tabular}
\end{center}

\vspace{-2pt}
(a) \ans{yes} \qquad (b) \ans{no --- it jumps at $x=3$}

\item \textbf{Reading intervals off a graph} (Lesson 2.4). Use graph (a) above. Name the stretch of
inputs on which the graph is \emph{constant}, and the stretch on which it is \emph{increasing}.

\ansline{Constant on $0 \le x \le 3$, that is $[0,3]$; increasing on $3 \le x \le 6$, that is $[3,6]$.}\\

\end{enumerate}
\end{notesbox}

\end{document}
